% Part: first-order-logic
% Chapter: sequent-calculus
% Section: rules-and-proofs

\documentclass[../../../include/open-logic-section]{subfiles}

\begin{document}

\iftag{FOL}
      {\olfileid[pt]{fol}{seq}{rul}}
      {\olfileid[pt]{pl}{seq}{rul}}

\olsection{Regras e \usetoken{P}{derivation}}

Para o que segue, sejam $\Gamma, \Delta, \Pi, \Lambda$ representando
sequências finitas de !!{sentence}s.

\begin{defn}[Sequente]
Um \emph{sequente} é uma expressão da forma
\[
\Gamma \Sequent \Delta
\]
em que $\Gamma$ e $\Delta$ são sequências finitas (possivelmente vazias) de
!!{sentence}s da linguagem $\Lang L$. $\Gamma$ é chamado de
\emph{antecedente}, enquanto $\Delta$ é o \emph{sucedente}.
\end{defn}

\begin{explain}
A ideia intuitiva por trás de um sequente é: se todas as !!{sentence}s no
antecedente valem, então ao menos uma das !!{sentence}s no
sucedente vale. Isto é, se $\Gamma = \tuple{!A_1, \dots, !A_m}$ e
$\Delta = \tuple{!B_1, \dots, !B_n}$, então $\Gamma \Sequent \Delta$
vale se, e somente se,
\[
(!A_1 \land \cdots \land !A_m) \lif (!B_1 \lor \cdots \lor
!B_n)
\]
vale. Há dois casos especiais: quando $\Gamma$ é vazio e quando
$\Delta$~é vazio. Quando $\Gamma$ é vazio, isto é, $m = 0$, $\quad
\Sequent \Delta$ vale se, e somente se, $!B_1 \lor \dots \lor !B_n$ vale. Quando
$\Delta$ é vazio, isto é, $n = 0$, $\Gamma \Sequent \quad$ vale se, e somente se,
$\lnot(!A_1 \land \dots \land !A_m)$ vale.  Dizemos que um sequente é válido
se, e somente se, a !!{sentence} correspondente é válida.
\end{explain}

Se $\Gamma$ é uma sequência de !!{sentence}s, escrevemos $\Gamma, !A$ para
o resultado de acrescentar $!A$ à extremidade direita de~$\Gamma$ (e $!A,
\Gamma$ para o resultado de acrescentar $!A$ à extremidade esquerda
de~$\Gamma$). Se $\Delta$ também é uma sequência de !!{sentence}s, então $\Gamma,
\Delta$ é a concatenação das duas sequências.

\begin{defn}[Sequente Inicial]
Um \emph{sequente inicial} é um sequente
\iftag{prvFalse,prvTrue}{de uma das seguintes formas:
  \begin{enumerate}
    \item $!A \Sequent !A$
    \tagitem{prvTrue}{$\quad \Sequent \ltrue$}{}
    \tagitem{prvFalse}{$\lfalse \Sequent \quad$}{}
  \end{enumerate}}
{da forma $!A \Sequent !A$} para qualquer !!{sentence} $!A$ na linguagem.
\end{defn}

!!^{derivation}s no cálculo de sequentes são certas árvores de
sequentes, em que os sequentes no topo são sequentes iniciais e, se um
sequente está abaixo de um ou dois outros sequentes, ele deve seguir
corretamente por uma regra de inferência.  As regras para $\Log{LK}$
dividem-se em dois tipos principais: regras \emph{lógicas} e
regras \emph{estruturais}.  As regras lógicas recebem o nome do !!{main operator}
da !!{sentence} que contém $!A$ e/ou $!B$ no
sequente inferior. Cada uma vem em duas versões, uma para inferir um
sequente com a !!{sentence} que contém o !!{operator} à esquerda,
e uma com a !!{sentence} à direita.

\end{document}
